\documentclass[11pt]{article}
\usepackage[rm2]{lecturenote}
\HandoutTitle{5}{The Discrete-Time Fourier Transform (DTFT)}{September 8, 2026}

\begin{document}
\maketitle

\HandoutMarkSection{1}{The Discrete-Time Fourier Transform (DTFT)}

\begin{defbox}[title={Def}]
The DTFT $X(e^{j\omega})$ of a sequence $x[n]$ is given by
\begin{equation}
X(e^{j\omega}) = \sum_{n=-\infty}^{\infty} x[n]\,e^{-j\omega n}.
\end{equation}
$\Rightarrow$ a complex-valued function of a real-valued variable $\omega$.
\begin{equation}
X(e^{j\omega}) = X_{\mathrm{re}}(e^{j\omega}) + j X_{\mathrm{im}}(e^{j\omega}).
\end{equation}
Both $X_{\mathrm{re}}(e^{j\omega})$ and $X_{\mathrm{im}}(e^{j\omega})$ are real-valued functions of $\omega$.
\end{defbox}

\begin{equation}
X(e^{j\omega}) = \underbrace{|X(e^{j\omega})|}_{\HandoutUnderbraceLabel{\text{magnitude}}}\,e^{j\HandoutHighlightYellow{$\theta(\omega)$}}, \qquad \theta(\omega) = \angle X(e^{j\omega})\ \text{(\textbf{phase})}.
\end{equation}

\begin{notebox}
Phase is not unique:
\begin{equation*}
  e^{j\theta(\omega)} = e^{j(\theta(\omega)+2\pi k)}, \quad k \in \Z.
\end{equation*}
\noindent Sometimes we force the \textbf{principal value},
$-\pi \le \theta(\omega) < +\pi$.
\end{notebox}

\begin{examplebox}[title={DTFT Examples}, handout mark=2]
\begin{itemize}
  \item $x[n] = \delta[n]$:
  \begin{align*}
    X(e^{j\omega}) &= \sum_{n=-\infty}^{\infty} \delta[n]\,e^{-j\omega n} = \delta[0]\,e^{-j\omega \cdot 0} = 1.
  \end{align*}
  \item $x[n] = \alpha^n u[n]$, $|\alpha|<1$:
  \begin{align*}
    X(e^{j\omega}) &= \sum_{n=-\infty}^{\infty} \alpha^n u[n]\,e^{-j\omega n}
      = \sum_{n=0}^{\infty} \alpha^n e^{-j\omega n} = \sum_{n=0}^{\infty} \left(\alpha e^{-j\omega}\right)^n \\
    &= \frac{1}{1-\alpha e^{-j\omega}} \quad \text{if } |\alpha e^{-j\omega}| < 1 \quad \left(|\alpha e^{-j\omega}| = |\alpha|\,\underbrace{|e^{-j\omega}|}_{\HandoutUnderbraceLabel{1}}\right).
  \end{align*}
\end{itemize}
\end{examplebox}

\begin{examplebox}[title={Example (Rectangular Window Sequence)}]
\begin{align*}
  x[n] &= \begin{cases} 1, & 0\le n \le N-1 \\ 0, & \text{else} \end{cases} \\
  &= u[n] - u[n-N].
\end{align*}%
\HandoutMarkBeforeAlign{3}%
\begin{align*}
  X(e^{j\omega}) &= \sum_{n=0}^{N-1} e^{-j\omega n} \\
  &= \frac{1-\left(e^{-j\omega}\right)^N}{1-e^{-j\omega}} \\
  &= \frac{e^{-j\omega\frac{N}{2}}\left(e^{+j\omega\frac{N}{2}}-e^{-j\omega\frac{N}{2}}\right)}{e^{-j\frac{\omega}{2}}\left(e^{+j\frac{\omega}{2}}-e^{-j\frac{\omega}{2}}\right)} \\
  &= \HandoutHighlightYellow{$\frac{\sin\!\left(\omega\frac{N}{2}\right)}{\sin\!\left(\frac{\omega}{2}\right)}$}\,e^{-j\omega\frac{N-1}{2}}.
\end{align*}
\end{examplebox}

\section{Dirichlet Kernel}

\begin{wrapfigure}[3]{r}{0.22\linewidth}
\vspace{-0.35\baselineskip}
\centering
\includegraphics[width=\linewidth,height=2.6cm,keepaspectratio]{Peter_Gustav_Lejeune_Dirichlet.jpg}\\[-0.1em]
{\scriptsize\itshape P.~G.~Lejeune Dirichlet}
\vspace{-0.35\baselineskip}
\end{wrapfigure}
In signal processing, the Dirichlet kernel is often called the periodic sinc function:
\[
\frac{\sin\!\left(\omega\frac{N}{2}\right)}{\sin\!\left(\frac{\omega}{2}\right)}.
\]
For small $\omega$, $\sin\!\left(\frac{\omega}{2}\right)\approx \frac{\omega}{2}$, so
\begin{equation}
\frac{\sin\!\left(\omega\frac{N}{2}\right)}{\sin\!\left(\frac{\omega}{2}\right)}
  \approx \frac{\sin\!\left(\omega\frac{N}{2}\right)}{\frac{\omega}{2}}
  = N\,\mathrm{sinc}\!\left(\omega\frac{N}{2}\right),
  \qquad \mathrm{sinc}(x) \triangleq \frac{\sin x}{x}.
\end{equation}
Zeros occur when $\omega\dfrac{N}{2} = k\pi$, $k\ne 0$, i.e., $\omega = \dfrac{2\pi k}{N}$.

\begin{center}
\adjustbox{max width=\linewidth,max height=0.24\textheight,center}{%
  \includegraphics{dirichlet_zeros.pdf}%
}\\[-0.35em]
{\scriptsize ``more zeros'' for larger $N$, faster oscillations.}
\end{center}

\noindent\HandoutMarkLeft{4}\textbf{But:} also periodic (``periodic sinc''):
\begin{align*}
\left|\frac{\sin\!\left(\omega\frac{N}{2}\right)}{\sin\!\left(\frac{\omega}{2}\right)}\right| \ \overset{?}{=}\ \left|\frac{\sin\!\left((\omega+2\pi)\frac{N}{2}\right)}{\sin\!\left(\frac{\omega+2\pi}{2}\right)}\right|
= \left|\frac{\sin\left(\omega\frac{N}{2}+\pi N\right)}{\sin\left(\frac{\omega}{2}+\pi\right)}\right|.
\end{align*}

\begin{center}
\adjustbox{max width=\linewidth,max height=0.22\textheight,center}{%
  \includegraphics{dirichlet_periodic.pdf}%
}\\[-0.35em]
{\scriptsize $|D(\omega)| = |D(\omega + 2\pi)|$}
\end{center}

\textbf{Dealing with $0/0$ when $\omega=0$:}
\begin{equation}
X(e^{j\omega})\Big|_{\omega=0} = \sum_{n=0}^{N-1} e^{-j\omega n}\Big|_{\omega=0} = N.
\end{equation}
Limits: $\displaystyle\lim_{\omega\to 0} \frac{f(\omega)}{g(\omega)}$. For small $\omega$, $\dfrac{\omega\frac{N}{2}}{\frac{\omega}{2}} = N$.

\HandoutMarkSection{5}{DTFT Properties}

\textbf{Periodicity:} The DTFT is $2\pi$-periodic.

\paragraph{Proof.}
\begin{align*}
X\!\left(e^{j(\omega_0+2\pi k)}\right) &= \sum_{n=-\infty}^{\infty} x[n]\,e^{-j(\omega_0+2\pi k)n} \\
&= \sum_{n=-\infty}^{\infty} x[n]\,e^{-j\omega_0 n}\,\underbrace{e^{-j2\pi k n}}_{\HandoutUnderbraceLabel{\left(e^{-j2\pi}\right)^{kn}=1}} \\
&= \sum_{n=-\infty}^{\infty} x[n]\,e^{-j\omega_0 n} = X(e^{j\omega_0}). \quad \blacksquare
\end{align*}

More properties: see \hyperref[tab:dtft-real-symmetry]{Table~3.1}, \hyperref[tab:dtft-complex-symmetry]{Table~3.2}, and \hyperref[tab:dtft-theorems]{Table~3.4} of Mitra.

\paragraph{Inverse.}
\begin{equation}
x[n] = \frac{1}{2\pi}\int_{-\pi}^{\pi} X(e^{j\omega})\,e^{+j\omega n}\,d\omega.
\end{equation}
\textbf{Proof.}
\begin{align*}
x[n] &= \frac{1}{2\pi}\int_{-\pi}^{\pi}\left(\sum_{\ell=-\infty}^{\infty} x[\ell]\,e^{-j\omega \ell}\right)e^{j\omega n}\,d\omega \\
&= \sum_{\ell=-\infty}^{\infty} x[\ell]\int_{-\pi}^{\pi}\frac{1}{2\pi}\,e^{j\omega(n-\ell)}\,d\omega \\
&= \sum_{\ell=-\infty}^{\infty} x[\ell]\,\frac{\sin\big(\pi(n-\ell)\big)}{\pi(n-\ell)}
= \sum_{\ell=-\infty}^{\infty} x[\ell]\,\delta[n-\ell] \\
&= x[n]. \quad \blacksquare
\end{align*}

\HandoutMarkSubsection{6}{Convergence}

\textbf{Note:} Exchanging the order of integration and summation (above) assumed convergence.

\begin{defbox}[title={Convergence}]
An infinite series of the form $X(e^{j\omega}) = \sum_{n=-\infty}^{\infty} x[n]\,e^{-j\omega n}$ may or may not converge. Let
\begin{equation}
X_K(e^{j\omega}) = \sum_{n=-K}^{K} x[n]\,e^{-j\omega n}, \qquad X(e^{j\omega}) = \lim_{K\to\infty} X_K(e^{j\omega}).
\end{equation}
\end{defbox}

\begin{defbox}[title={Def}]
We say that $X_K(e^{j\omega})$ converges \textbf{uniformly} to $X(e^{j\omega})$ if for every $\varepsilon>0$, there exists an integer $K$ such that for all integers $L\ge K$,
\begin{equation}
\big|X(e^{j\omega}) - X_L(e^{j\omega})\big| < \varepsilon \quad \text{for all } \omega.
\end{equation}
\end{defbox}

\begin{factbox}[title={Property}]
The DTFT will converge uniformly if $x[n]$ is an \textbf{absolutely summable} sequence,
\begin{equation}
\sum_{n=-\infty}^{\infty} |x[n]| < \infty.
\end{equation}
\end{factbox}

\HandoutMarkLeft{7}

\begin{defbox}[title={Def}]
We say that $X_K(e^{j\omega})$ converges in the \textbf{mean-square} to $X(e^{j\omega})$ if
\begin{equation}
\lim_{K\to\infty} \int_{-\pi}^{\pi} \big|X(e^{j\omega}) - X_K(e^{j\omega})\big|^2\,d\omega = 0.
\end{equation}
\end{defbox}

\begin{factbox}[title={Property}]
The DTFT will converge in the mean-square if the sequence $x[n]$ is \textbf{square-summable}:
\begin{equation}
\sum_{n=-\infty}^{\infty} |x[n]|^2 < \infty.
\end{equation}
\end{factbox}

\begin{factbox}[title={Property (Nesting of the $\ell_p(\mathbb{Z})$ spaces)}]
\begin{itemize}
  \item Every absolutely summable sequence is also square summable.
  \item Every square-summable sequence is bounded.
\end{itemize}
\begin{equation}
\ell_1(\mathbb{Z}) \subset \ell_2(\mathbb{Z}) \subset \ell_\infty(\mathbb{Z}).
\end{equation}
\begin{handouttikz}
  \definecolor{swissred}{RGB}{227,6,19}
  \tikzset{
    lp/lbl/.style={
      align=left,
      text=ee483ink,
      execute at begin node={\footnotesize\sffamily}
    },
    lp/ann/.style={
      align=left,
      text=ee483ink,
      execute at begin node={\large\sffamily}
    },
    lp/lead/.style={
      draw=ee483ink,
      line width=0.65pt,
      line cap=rect,
      -{Stealth[length=2.8mm, width=2.1mm, inset=0pt]}
    },
    lp/leadaccent/.style={
      draw=swissred,
      line width=0.75pt,
      line cap=rect,
      -{Stealth[length=2.8mm, width=2.1mm, inset=0pt]}
    }
  }

  \draw[ee483ink, line width=0.35pt] (-6.2,3.85) -- (6.2,3.85);

  \draw[ee483ink, line width=0.5pt, fill=white]
    (-6,-3.5) rectangle (6,3.5);

  \draw[ee483ink, line width=0.75pt]
    (0,-0.1) ellipse (5cm and 2.7cm);

  \filldraw[
    fill=ee483ink!4,
    draw=ee483ink,
    line width=0.9pt
  ]
  (-0.8,0) ellipse (3.3cm and 1.8cm);

  \filldraw[
    fill=swissred!10,
    draw=swissred,
    line width=1.1pt
  ]
  (-0.8,0) ellipse (1.25cm and 0.8cm);

  \node[lp/lbl, anchor=north east] at (5.85, 3.25)
    {Space of all sequences};

  \node[lp/ann, anchor=west] (lblinf) at (6.35, -2.05)
    {Bounded sequences\\[1pt]$\ell_\infty(\mathbb{Z})$};
  \draw[lp/lead] (lblinf.west) ++(-0.05, 0.32) -- (3.85, -0.92);

  \node[lp/ann, anchor=east] (lbltwo) at (-6.35, 1.55)
    {Square-summable\\[1pt]$\ell_2(\mathbb{Z})$};
  \draw[lp/lead] (lbltwo.east) ++(0.05, -0.18) -- (-3.15, 0.82);

  \node[
    lp/ann,
    align=center,
    text=swissred
  ] (lblone) at (-0.8,-2.35)
    {Absolutely summable\\[1pt]$\ell_1(\mathbb{Z})$};
  \draw[lp/leadaccent] (lblone.north) ++(0, 0.08) -- (-0.8,-1.05);
\end{handouttikz}
\end{factbox}

\HandoutMarkSubsection{8}{An Example, and Generalized DTFTs}

\begin{examplebox}[title={Example}]
The sequence
\begin{equation}
x[n] = \begin{cases} 1/n, & n \ge 1 \\ 0, & n \le 0 \end{cases}
\end{equation}
is \textbf{not} absolutely summable, but \textbf{is} square summable (famous result by Euler):
\begin{equation}
\sum_{n=1}^{\infty} \left(\frac{1}{n}\right)^2 = \frac{\pi^2}{6}.
\end{equation}
\end{examplebox}

We can sometimes define DTFTs in cases where the DTFT does not actually converge (but still exists in a generalized sense).

\begin{examplebox}[title={Example}]
$x[n] = e^{j\omega_0 n}$:
\begin{equation}
X(e^{j\omega}) = \sum_{n=-\infty}^{\infty} e^{j(\omega_0-\omega)n} \qquad \text{does not converge.}
\end{equation}
But $x[n]$ still has a useful (generalized) DTFT:
\begin{equation}
X(e^{j\omega}) = \sum_{k=-\infty}^{\infty} 2\pi\,\delta(\omega - \omega_0 + 2\pi k).
\end{equation}
\end{examplebox}

\HandoutMarkLeft{9}

\begin{equation}
X(e^{j\omega}) = 2\pi\,
\underbrace{\delta(\omega-\omega_0)}_{\HandoutUnderbraceLabel{\text{Dirac delta function}}},
\quad -\pi\le\omega\le\pi,\ -\pi\le\omega_0\le\pi
\end{equation}
\begin{align*}
x[n] &= \frac{1}{2\pi}\int_{-\pi}^{\pi} 2\pi\,\delta(\omega-\omega_0)\,e^{j\omega n}\,d\omega = e^{j\omega_0 n}.
\end{align*}

\section*{Readings}
\begin{itemize}
  \item Sanjit K. Mitra, \textit{Digital Signal Processing: A Computer-Based Approach}, 4th ed., Ch.\ 4.8
  \item Monson H. Hayes, \textit{Schaum's Outline of Digital Signal Processing}, 2nd ed., Ch.\ 2.7
\end{itemize}

\HandoutAppendixListStart
  \HandoutAppendixListItem{app:mitra-tables}{A.\ Mitra Tables 3.1, 3.2, and 3.4}
  \HandoutAppendixListItem{app:scan5}{B.\ Handwritten Lecture Note Scan (September 8, 2026)}
\HandoutAppendixListEnd

\HandoutAppendixSection{app:mitra-tables}{Mitra Tables 3.1, 3.2, and 3.4}
\ifhandoutclean\else
% Mitra 4th ed.: Tables 3.1, 3.2, and 3.4 (DTFT properties appendix).

\begin{table}[H]
\centering
\caption{Symmetry relations of the discrete-time Fourier transform of a real sequence (Mitra Table~3.1).}
\label{tab:dtft-real-symmetry}

\begin{tabular}{cc}
\toprule
\textbf{Sequence}
&
\textbf{Discrete-Time Fourier Transform}
\\
\midrule

$x[n]$
&
$X(e^{j\omega})
 = X_{\mathrm{re}}(e^{j\omega})
 + jX_{\mathrm{im}}(e^{j\omega})$
\\
\midrule

$x_{\mathrm{ev}}[n]$
&
$X_{\mathrm{re}}(e^{j\omega})$
\\

$x_{\mathrm{od}}[n]$
&
$jX_{\mathrm{im}}(e^{j\omega})$
\\
\midrule

\multirow{5}{*}{Symmetry relations}
&
$X(e^{j\omega})=X^{*}(e^{-j\omega})$
\\[1.2em]

&
$X_{\mathrm{re}}(e^{j\omega})
 =X_{\mathrm{re}}(e^{-j\omega})$
\\[1.2em]

&
$X_{\mathrm{im}}(e^{j\omega})
 =-X_{\mathrm{im}}(e^{-j\omega})$
\\[1.2em]

&
$\lvert X(e^{j\omega})\rvert
 =\lvert X(e^{-j\omega})\rvert$
\\[1.2em]

&
$\arg\!\left\{X(e^{j\omega})\right\}
 =-\arg\!\left\{X(e^{-j\omega})\right\}$
\\

\bottomrule
\end{tabular}

\medskip

\parbox{0.9\textwidth}{
\textit{Note:}
$x_{\mathrm{ev}}[n]$ and $x_{\mathrm{od}}[n]$ denote the even and
odd parts of $x[n]$, respectively.
}
\end{table}

\begin{table}[H]
\centering
\caption{Symmetry relations of the discrete-time Fourier transform of a complex sequence (Mitra Table~3.2).}
\label{tab:dtft-complex-symmetry}

\begin{tabular}{cc}
\toprule
\textbf{Sequence}
&
\textbf{Discrete-Time Fourier Transform}
\\
\midrule

$x[n]$
&
$X(e^{j\omega})$
\\
\midrule

$x[-n]$
&
$X(e^{-j\omega})$
\\[0.8em]

$x^{*}[-n]$
&
$X^{*}(e^{j\omega})$
\\[0.8em]

$x_{\mathrm{re}}[n]$
&
$X_{\mathrm{cs}}(e^{j\omega})
 =\dfrac{1}{2}
 \left\{
 X(e^{j\omega})+X^{*}(e^{-j\omega})
 \right\}$
\\[0.8em]

$jx_{\mathrm{im}}[n]$
&
$X_{\mathrm{ca}}(e^{j\omega})
 =\dfrac{1}{2}
 \left\{
 X(e^{j\omega})-X^{*}(e^{-j\omega})
 \right\}$
\\[0.8em]

$x_{\mathrm{cs}}[n]$
&
$X_{\mathrm{re}}(e^{j\omega})$
\\[0.8em]

$x_{\mathrm{ca}}[n]$
&
$jX_{\mathrm{im}}(e^{j\omega})$
\\

\bottomrule
\end{tabular}

\medskip

\parbox{0.9\textwidth}{
\textit{Note:}
$X_{\mathrm{cs}}(e^{j\omega})$ and
$X_{\mathrm{ca}}(e^{j\omega})$ are the conjugate-symmetric and
conjugate-antisymmetric parts of $X(e^{j\omega})$, respectively.
Likewise, $x_{\mathrm{cs}}[n]$ and $x_{\mathrm{ca}}[n]$ are the
conjugate-symmetric and conjugate-antisymmetric parts of $x[n]$,
respectively.
}
\end{table}

\begin{table}[H]
\centering
\caption{Discrete-time Fourier transform theorems (Mitra Table~3.4).}
\label{tab:dtft-theorems}

\small
\begin{tabular}{@{}>{\raggedright\arraybackslash}p{0.17\linewidth}@{}>{\raggedright\arraybackslash}p{0.28\linewidth}@{}>{\raggedright\arraybackslash}p{0.47\linewidth}@{}}
\toprule
\textbf{Theorem} & \textbf{Sequence} & \textbf{Discrete-Time Fourier Transform} \\
\midrule
& $g[n]$ & $G(e^{j\omega})$ \\
& $h[n]$ & $H(e^{j\omega})$ \\
\midrule
Linearity &
$\alpha g[n]+\beta h[n]$ &
$\alpha G(e^{j\omega})+\beta H(e^{j\omega})$ \\[0.6em]
Time-reversal &
$g[-n]$ &
$G(e^{-j\omega})$ \\[0.6em]
Time-shifting &
$g[n-n_0]$ &
$e^{-j\omega n_0}G(e^{j\omega})$ \\[0.6em]
Frequency-shifting &
$e^{j\omega_0 n}g[n]$ &
$G\!\left(e^{j(\omega-\omega_0)}\right)$ \\[0.6em]
Differentiation-in-frequency &
$ng[n]$ &
$j\dfrac{dG(e^{j\omega})}{d\omega}$ \\[0.8em]
Convolution &
$g[n]\circledast h[n]$ &
$G(e^{j\omega})H(e^{j\omega})$ \\[0.6em]
Modulation &
$g[n]h[n]$ &
$\dfrac{1}{2\pi}\displaystyle\int_{-\pi}^{\pi}G(e^{j\theta})H(e^{j(\omega-\theta)})\,d\theta$ \\[0.8em]
Parseval's Theorem &
\multicolumn{2}{@{}l@{}}{%
$\displaystyle\sum_{n=-\infty}^{\infty}g[n]h^*[n]
=\dfrac{1}{2\pi}\displaystyle\int_{-\pi}^{\pi}G(e^{j\omega})H^*(e^{j\omega})\,d\omega$}%
\\
\bottomrule
\end{tabular}
\end{table}

\clearpage
\fi

\HandoutAppendixSection{app:scan5}{Handwritten Lecture Note Scan}
\HandoutIncludeHandoutScan{lecture5_0908.pdf}

\end{document}
